\scanpage{229}
\mspage{762}

point $1$ of $\mathbb U$, so that the universal covering of
$\mathbb S_0$ pointed over $e_0$
is identified with the exponential \unsure{covering}
\[
\tag{11}
\diagram{\begin{tikzcd}[ampersand replacement=\&,column sep=large,row sep=small]
\mathbb R \arrow[r] \& \mathbb U \&
(\xrightarrow{\sim}\mathbb S_0)\\
t \arrow[r,mapsto] \& \exp 2i\pi t \&
(\mapsto \mathbb R^+\exp 2i\pi t.
\end{tikzcd}}
\]

Thus, a lifting of $u_{\unsure{\mathbb S^1}}$
\inserted{to $\widetilde u$} amounts to the choice of $t\in\mathbb R$ such that
\[
\tag{12}
\exp 2i\pi t=\frac{u(1)}{|u(1)|}\qquad
\deleted{\unsure{\text{is written}}}
\]
i.e. to the choice of an argument $t$ of $u(1)$.

We apply the construction of $\widetilde G$ of $1^\circ)$
to the case of $G=\mathrm{Gl}(2,\mathbb R)^{\pm1}$,
we find an extension structure
\[
\tag{13}
\diagram{\begin{tikzcd}[ampersand replacement=\&,column sep=1.4em]
1\arrow[r]\&\mathbb Z\arrow[r]\&
\mathrm{Gl}(2,\mathbb R)^{\pm1\,\sim}\arrow[r]\&
\mathrm{Gl}(2,\mathbb R)^{\pm1}\arrow[r]\&1.
\end{tikzcd}}
\]

\deleted{\unsure{Propo}}\inserted{which induces}
over the identity component $\mathrm{Sl}(2,\mathbb R)$
of $\mathrm{Gl}(2,\mathbb R)^{\pm1}$ an extension structure
\[
\tag{14}
\diagram{\begin{tikzcd}[ampersand replacement=\&,column sep=1.5em]
1\arrow[r]\&\mathbb Z\arrow[r]\&
\mathrm{Sl}(2,\mathbb R)^{\sim}\arrow[r]\&
\mathrm{Sl}(2,\mathbb R)\arrow[r]\&1.
\end{tikzcd}}
\]

\begin{tabular}{|p{0.9\linewidth}}
\uline{Proposition} The group $\mathrm{Sl}(2,\mathbb R)^{\sim}$ is
$1$-connected, and even \deleted{\unsure{contra}}
homeomorphic to $\mathbb R^3$.
\end{tabular}

\margin{228}

\begin{center}
228\\
\texttt{https://grothendieck.umontpellier.fr}\\
Université de Montpellier, UMR CNRS 5149, Case courrier 051, Place Eugène Bataillon\\
34095 Montpellier cedex 5 - France
\end{center}

\scanpage{230}

\uline{Proof} Since the inclusion
\[
\tag{15}
\diagram{\begin{tikzcd}[ampersand replacement=\&,column sep=large]
\mathbb U\arrow[r,hook]\&\mathrm{Sl}(2,\mathbb R)
\end{tikzcd}}
\]
is a homotopy equivalence,
\deleted{\unsure{which}} is identified \inserted{topologically}
with the inclusion
$\diagram{\begin{tikzcd}[ampersand replacement=\&,column sep=large]
\mathbb U\arrow[r]\&\mathbb U\times\mathbb R^2
\end{tikzcd}}$,
we are reduced to proving that
\deleted{\unsure{over} $\mathbb U$}
\inserted{the inverse image $\mathbb U^{\sim}$}
of $\mathbb U$ in $\mathrm{Sl}(2,\mathbb R)^{\sim}$
is identified with the universal covering
$\deleted{\unsure{\mathbb R}}\,\mathbb R$ of $\mathbb U$
(pointed over $1$ by $0\ldots$)
Now if $u\in\mathbb U$, the \deleted{\unsure{last}}
choice of a $\widetilde u$ over $u$, i.e. of a lifting of
$\dfrac{u(1)}{|u(1)|}$ to $\mathbb R$ by formula (12),
is identified (since $\dfrac{u(1)}{|u(1)|}=u(1)=u$)
with the choice of a $t\in\mathbb R$ such that $\exp 2i\pi t\in u$, O.K.

\uline{Corollary 1} The \unsure{inverse} \unsure{subgroup} in
$\mathrm{Sl}(\unsure{2},\mathbb R)^{\sim}$ of a
finite cyclic subgroup \inserted{of order $n$} is $\simeq\mathbb Z$
($\exists!$ generator $\overline\omega$ of it,
such that we have $\overline\omega^{\,n}=\omega'_0$,
where $\omega'_0$ is the \deleted{\unsure{canonical}} generator
of the $\operatorname{Ker}(\mathrm{Sl}(2,\mathbb R)^{\sim}
\to\mathrm{Sl}(2,\mathbb R))$).

\margin{Applying this to the center
$\{\pm1\}$ of $\mathrm{Sl}(2,\mathbb R)$,
we find a canonical element $\omega_0\in\mathrm{Sl}(2,\mathbb R)^{\sim}$
with $\omega_0^2=\omega'_0$.}

\uline{Corollary 2} The action of a $g\in\mathrm{Gl}(2,\mathbb R)^{-1}$
on $\deleted{\unsure{\overline\omega}}\,\omega'_0$
(or on the center $\omega_0^{\mathbb Z}$ of
$\mathrm{Sl}(2,\mathbb R)^{\sim}$)
is via $\omega'_0\mapsto{\omega'_0}^{-1}$
(resp. $\omega_0\mapsto\omega_0^{-1}$) $\ldots$

\margin{229}

\begin{center}
229\\
\texttt{https://grothendieck.umontpellier.fr}\\
Université de Montpellier, UMR CNRS 5149, Case courrier 051, Place Eugène Bataillon\\
34095 Montpellier cedex 5 - France
\end{center}

\scanpage{231}
\mspage{763}

$3^\circ)$ \uline{The group} $\mathrm{Gl}(2,\mathbb Z)^{\sim}$.

Let $\mathrm{Gl}(2,\mathbb Z)^{\sim}$ be the inverse image of
$\mathrm{Gl}(2,\mathbb Z)$ in $\mathrm{Gl}(2,\mathbb R)^{\pm1\,\sim}$,
so that we have an extension structure
\[
\tag{16}
\diagram{\begin{tikzcd}[ampersand replacement=\&,column sep=1.2em,row sep=small]
1\arrow[r]\&\mathbb Z\arrow[r]\&\mathrm{Gl}(2,\mathbb Z)^{\sim}\arrow[r]\&
\mathrm{Gl}(2,\mathbb Z)\arrow[r]\&1.\\
\&\deleted{\unsure{n}}\arrow[u,mapsto]\arrow[r,mapsto]\&
\deleted{\unsure{\omega'_0}}\\
\&n\arrow[r,mapsto]\&\unsure{\omega_0^n|{\omega'_0}^{\,n}}
\end{tikzcd}}
\]

I would like to describe $\mathrm{Gl}(2,\mathbb Z)^{\sim}$ by generators
and relations. For this, I note this:

\begin{tabular}{|p{0.9\linewidth}}
\uline{Proposition} The inverse image
\inserted{in $\mathrm{Gl}(2,\mathbb Z)^{\sim}$}
of the cyclic subgroups $L_\rho\simeq\mathbb Z/6\mathbb Z$
and $L_\sigma\simeq\mathbb Z/4\mathbb Z$ generated by
$\rho$ and $\sigma$ are isomorphic to $\mathbb Z$.
\end{tabular}

This follows from the

\uline{Corollary} (to the proposition p. 762): For every
finite cyclic subgroup $\unsure{\mathfrak Z|\mathfrak z}\subset
\mathrm{Sl}(2,\mathbb R)$, the inverse image of
$\unsure{\mathfrak Z|\mathfrak z}$ in
$\mathrm{Sl}(2,\mathbb R)^{\sim}$ is isomorphic to $\mathbb Z$.

Indeed, by conjugation \unsure{we may assume} that
$\unsure{\mathfrak Z|\mathfrak z}\subset\mathbb U$,
i.e. $\unsure{\mathfrak Z|\mathfrak z}=\mu_n(\mathbb C)$,
but then $\unsure{\widetilde{\mathfrak Z}|\widetilde{\mathfrak z}}$
is identified with the set of $t\in\mathbb R$ with
$(\exp 2i\pi t)^n=1$ i.e. $\exp 2i\pi(nt)=1$ i.e.

\margin{230}

\begin{center}
230\\
\texttt{https://grothendieck.umontpellier.fr}\\
Université de Montpellier, UMR CNRS 5149, Case courrier 051, Place Eugène Bataillon\\
34095 Montpellier cedex 5 - France
\end{center}

\scanpage{232}

\[
t\in\frac1n\mathbb Z,
\]
the group of these $t$ is $\simeq\mathbb Z$, q.e.d.

\uline{Corollary} There exist
$\widetilde\rho\in\mathrm{Sl}(2,\mathbb Z)^{\sim}$ and
$\widetilde\sigma\in\mathrm{Sl}(2,\mathbb Z)^{\sim}$,
over $\rho$ and $\sigma$ respectively, such that we have
\[
\tag{17}
\widetilde\rho^{\,3}=\widetilde\sigma^{-2}
\]
\margin{for a more \unsure{pleasant} statement, see (20) below $\ldots$}
The set of these systems $(\widetilde\rho,\widetilde\sigma)$
is a torsor under $\mathbb Z$, by the operation
\[
\tag{18}
\diagram{\begin{tikzcd}[ampersand replacement=\&,column sep=large]
(\widetilde\rho,\widetilde\sigma)\arrow[r,mapsto]\&
(\widetilde\rho\,{\omega'_0}^{\,2n},
\widetilde\sigma\,{\omega'_0}^{\,3n}),
\end{tikzcd}}
\]
where $\omega'_0$ is the canonical generator of
$\diagram{\begin{tikzcd}[ampersand replacement=\&,]
\mathrm{Sl}(2,\mathbb Z)^{\sim}\arrow[r]\&
\mathrm{Sl}(2,\mathbb Z)
\end{tikzcd}}$.

We can also say that the inverse image of the element
$\deleted{\unsure{=-\mathrm{id}}}$
\inserted{of the center $\{\pm1\}$} of $\mathrm{Sl}(2,\mathbb Z)$
\deleted{\unsure{the unique}}
\deleted{(the central element \unsure{distinct from|$\ne$}
$+1$ of $\mathrm{Sl}(2,\mathbb Z)$)}
\deleted{is a torsor under $\mathbb Z$, and the choice}
\inserted{is a group isomorphic to $\mathbb Z$, with a}
unique generator $\omega_0$ satisfying
\deleted{\unsure{of a} lifting}
\[
\tag{19}
\begin{gathered}
\omega_0^2=\omega'_0
\quad(\inserted{\text{the}}\ \text{canonical generator of}\\
\diagram{\begin{tikzcd}[ampersand replacement=\&,]
\mathrm{Sl}(2,\mathbb Z)^{\sim}\arrow[r]\&
\mathrm{Sl}(2,\mathbb Z)
\end{tikzcd}})
\end{gathered}
\]

\margin{it \unsure{is} better to say: $\exists!$ generator
$\widetilde\rho$ \inserted{(resp. $\widetilde\sigma$)}
of the inverse image $\unsure{\widetilde L_\rho|L_\rho}$ of $L_\rho$
(resp. $L_\sigma$) in $\mathrm{Sl}(2,\mathbb Z)^{\sim}$,
over $\rho$ (resp. $\sigma$), and we then have relations (20)}

and this being so, it is possible to lift $\rho$ to $\widetilde\rho$,
\inserted{$\sigma$ to $\widetilde\sigma$} satisfying
\[
\tag{20}
\widetilde\rho^{\,3}=\omega_0^{-1},\qquad
\widetilde\sigma^{\,2}=\omega_0
\]
\margin{\uline{Do the verification}}
And this time, $\widetilde\rho$ and $\widetilde\sigma$
are uniquely deter-

\margin{231}

\begin{center}
231\\
\texttt{https://grothendieck.umontpellier.fr}\\
Université de Montpellier, UMR CNRS 5149, Case courrier 051, Place Eugène Bataillon\\
34095 Montpellier cedex 5 - France
\end{center}

\scanpage{233}
\mspage{764}

mined. This follows \deleted{\unsure{in a way}}
\inserted{\deleted{\unsure{from the}}}
at once from the preceding proposition.

\unsure{i)} Let us consider the inverse image of $\{1,\tau\}\subset\mathrm{Gl}(2,\mathbb Z)$
in $\mathrm{Gl}(2,\mathbb Z)^{\sim}$,
\deleted{\unsure{a group}} extension of
$\{1,\tau\}\simeq\mathbb Z/2\mathbb Z$ by $\mathbb Z$
--- I say that this extension is \uline{split}.

\margin{\uline{N.B.}
$\tau=\tau_\infty=\begin{pmatrix}-1&0\\0&1\end{pmatrix}$
i.e. $\tau(x+iy)=\unsure{-}x+iy$
(reflection with respect to the $y$-axis)}

Indeed, the action of $\tau$ on $\mathbb S_0$
lifts to the action of $\widetilde\tau$ \deleted{of}
of order $2$ on
$\deleted{\unsure{\mathbb R}}\,
\widetilde{\mathbb S}_0\simeq\mathbb R$,
which is given by
\[
\tag{21}
\widetilde\tau t=\frac12-t
\]
for we have for $z\in\mathbb U$, $z=\exp 2i\pi t$
\[
\tag{22}
\begin{aligned}
\tau z&=-\overline z=-\exp 2i\pi t
=\exp(\pi it-2i\pi t)\\
&=\exp\unsure{\{}\,2i\pi
\underbrace{\left(\frac12-t\right)}_{\widetilde\tau(t)}.
\end{aligned}
\]

\deleted{I} It is perhaps more convenient to work with
$\tau'=\tau'_\infty=\begin{pmatrix}0&1\\1&0\end{pmatrix}$,
which \deleted{is} exchanges $e_0=1$ and $e_1=i$,
i.e. the reflection \unsure{with respect|by} to the main diagonal.
This automorphism stabilizes $\mathbb U$, and fixes the point
$\exp\dfrac{\pi i}{\unsure{4|8}}=\dfrac{\sqrt2}{2}(i+1)\unsure{)}$,
it lifts to
\[
\tag{23}
\widetilde{\tau'}(t)=2i\pi\left(\frac1{\unsure{8|4}}-t\right).
\]
Let us note that \deleted{$\unsure{\widetilde{\tau'}|\tau'}$}
the action of $\widetilde{\tau'}$ (like that

\margin{232}

\begin{center}
232\\
\texttt{https://grothendieck.umontpellier.fr}\\
Université de Montpellier, UMR CNRS 5149, Case courrier 051, Place Eugène Bataillon\\
34095 Montpellier cedex 5 - France
\end{center}

\scanpage{234}

of $\widetilde\tau$) on $\omega_0$ is the symmetry $\omega_0\mapsto\omega_0^{-1}$ (cf. end of the preceding \S\ ). I say that we have
\[
\begin{array}{ll}
(24)&\widetilde\tau'(\widetilde\rho)=\widetilde\rho^{-1},\qquad
\widetilde\tau'(\widetilde\sigma)=\widetilde\sigma^{-1}.
\end{array}
\]
The formulas are \unsure{also|thus} written
\[
\widetilde\tau'(\widetilde\rho^{-1})=\widetilde\rho,\qquad
\widetilde\tau'(\widetilde\sigma^{-1})=\widetilde\sigma
\]
they are valid \deleted{\unsure{and}} in $\unsure{\mathrm{Sl}}(2,\mathbb Z)$, so to see that they are valid in $\mathrm{Sl}(2,\mathbb Z)^{\sim}$, it suffices to show that the left-hand \unsure{sides} satisfy the \unsure{characteristic} relations (20), which is written
\[
\widetilde\tau'\bigl(\underbrace{\widetilde\rho^{-3}}_{\omega_0^{\deleted{\unsure{-1}}}}\bigr)=\omega_0^{-1},\qquad
\widetilde\tau'\bigl(\underbrace{\widetilde\sigma^{-2}}_{\omega_0^{-1}}\bigr)=\omega_0
\]
i.e. $\widetilde\tau'(\omega_0)=\omega_0^{-1}$, OK.

\begin{tabular}{|l}
\uline{Theorem.} The group $\mathrm{Sl}(2,\mathbb Z)^{\sim}$ admits\\
the presentation (with two generators $\widetilde\rho,\widetilde\sigma$)\\[3pt]
$(25)\qquad\boxed{\widetilde\rho^3\cdot\widetilde\sigma^2=1}$
\end{tabular}

\uline{Proof} Let $\unsure{\Pi|\pi}$ be the group \deleted{\unsure{generated}} \inserted{defined} by this presentation, and let us consider the homomorphism
\diagram{
\begin{tikzcd}[ampersand replacement=\&,column sep=large,row sep=tiny]
\unsure{\Pi|\pi}\arrow[r] \& \mathrm{Sl}(2,\mathbb Z)/\pm1=\{\rho,\sigma\mid\rho^3=\sigma^2=1\}\\
(\widetilde\rho,\widetilde\sigma)\arrow[r,mapsto] \& (\rho,\sigma)\ \mathrm{mod}\{\pm1\}
\end{tikzcd}
\quad(26)
}

\margin{233}

\begin{center}
233\\
\texttt{https://grothendieck.umontpellier.fr}\\
Université de Montpellier, UMR CNRS 5149, Case courrier 051, Place Eugène Bataillon\\
34095 Montpellier cedex 5 - France
\end{center}

\scanpage{235}
\mspage{765}

the group $\mathrm{Sl}(2,\mathbb Z)/\pm1$ is identified with the quotient of $\unsure{\pi|\Pi}$ by the normal subgroup generated by the element $\uline{\omega_0}\overset{\mathrm{dfn}}{=}\widetilde\rho^{-3}=\widetilde\sigma^2$.
But this element being \unsure{central}, the group generated by $\uline{\omega_0}$ is already normal.
thus we have an exact sequence
\diagram{
$(27)\quad$
\begin{tikzcd}[ampersand replacement=\&,column sep=large,row sep=tiny]
n\arrow[r,mapsto] \& \uline{\omega_0}^{n} \& \& \\
\mathbb Z\arrow[r] \& \unsure{\Pi|\pi}\arrow[r] \& \mathrm{Sl}(2,\mathbb Z)/\pm1\arrow[r] \& 1
\end{tikzcd}
}
which maps into the exact sequence
\diagram{
\deleted{$\unsure{(28)|(27)}$}\quad
\begin{tikzcd}[ampersand replacement=\&,column sep=large,row sep=tiny]
1\arrow[r] \& \mathbb Z\arrow[r] \& \mathrm{Sl}(2,\mathbb Z)^{\sim}\arrow[r] \& \mathrm{Sl}(2,\mathbb Z)/\pm1\arrow[r] \&1\\
\& n\arrow[r,mapsto] \&\omega_0^{n}\&\&
\end{tikzcd}
}
inducing \inserted{identity} isomorphisms on the end factors --- thus by the lemma \unsure{of the|of} $5$
$\unsure{\pi|\Pi}\longrightarrow\mathrm{Sl}(2,\mathbb Z)^{\sim}$ is an isomorphism, q.e.d.

\begin{tabular}{|l}
\uline{Corollary 1} The group $\mathrm{Gl}(2,\mathbb Z)^{\sim}$ admits\\
the presentation (with three generators $\widetilde\rho,\widetilde\sigma,\widetilde\tau'$)
\end{tabular}
\[
\begin{array}{ll}
(28)&\boxed{
\widetilde\tau'{}^{2}=\deleted{\unsure{1,|\omega_0}}\,
\widetilde\rho^3\cdot\widetilde\sigma^2=1,\quad
\widetilde\tau'(\widetilde\rho)=\widetilde\rho^{-1},\quad
\widetilde\tau'(\widetilde\sigma)=\widetilde\sigma^{-1}}
\end{array}
\]

I am going to \unsure{try} to present the \deleted{\unsure{element}} \inserted{group}
$\mathrm{Gl}(2,\mathbb Z)^{\sim}$ \unsure{with|in} three involutive generators,
\deleted{\unsure{at the of}} $\widetilde r_0,\widetilde r_1,\widetilde r_\infty$ over the elements
$r_0,\unsure{r_1|\sigma_0},r_\infty$ of $\mathrm{Gl}(2,\mathbb Z)^{-}$, by \unsure{drawing inspiration} \unsure{from the formulary|from the formularies} of $\mathrm{Sl}(2,\mathbb Z)$:

\margin{234}

\begin{center}
234\\
\texttt{https://grothendieck.umontpellier.fr}\\
Université de Montpellier, UMR CNRS 5149, Case courrier 051, Place Eugène Bataillon\\
34095 Montpellier cedex 5 - France
\end{center}